% ed.: Consolidated from the four Legendre--Rvachev reports listed in
% Appendix~\ref{sec:provenance-poly}.  The symbols Q^- and R deliberately
% distinguish deconvolution from moment smoothing.

\section{Legendre synthesis, spectral closure, and self-reconstruction}
\label{sec:legendre-closure}

The Legendre reports compose two exact transforms: Fourier--Legendre analysis
of the up-function and finite synthesis of each polynomial mode by up-atoms.
Their principal notational ambiguity is removed once and for all by the box
\begin{equation}
 \boxed{
 \begin{gathered}
 Q_n^-:=\MM_u(D)^{-1}P_n,\qquad
 R_n:=\MM_u(D)P_n,\\
 X_-:=\MM_u(D)^{-1}x\MM_u(D),\qquad
 X_+:=\MM_u(D)x\MM_u(D)^{-1}.
 \end{gathered}}
 \label{eq:leg-notation}
\end{equation}
Thus $Q_n^-$ is the deconvolved polynomial whose samples are synthesis
coefficients, whereas $R_n$ is the distinct moment-smoothed polynomial used
in the Appell connection of \cref{sec:leg-smoothed}.  If
$K(z)=\log\MM_u(z)$, the commutator identity
$g(D)x=xg(D)+g'(D)$ gives
\begin{equation}
 X_-=x-K'(D),\qquad X_+=x+K'(D).
 \label{eq:leg-Xpm}
\end{equation}
We also write $\mathscr D_u:=\MM_u(D)^{-1}$ for the deconvolution operator.

Use the classical Legendre normalization $P_n(1)=1$, for which
\begin{equation}
 h_n=\int_{-1}^{1}P_n(x)^2\Differential x=\frac2{2n+1}
 \label{eq:leg-hn}
\end{equation}
and define the even Fourier--Legendre coefficients
\begin{equation}
 a_n:=\frac{4n+1}{2}\int_{-1}^{1}u(x)P_{2n}(x)\Differential x.
 \label{eq:leg-an}
\end{equation}
The primary exposition's section on the Fourier--Legendre expansion of
$\RvachevUp$ in \path{Fabius_Function_and_Rvachev_Up.tex} proves
\begin{equation}
 u(x)=\sum_{n=0}^{\infty}a_nP_{2n}(x)
 \label{eq:leg-canonical-series}
\end{equation}
absolutely and uniformly on $[-1,1]$, with rational $a_n$; the exact Lean
anchors are \path{summable_abs_rvachevLegendreCoefficient},
\path{hasSum_rvachevLegendreSeries}, and
\path{hasSum_rvachevLegendreSeries_uniform} in
\path{FabiusLegendreSeries.lean}.

\subsection{Finite mode synthesis and the legal block loop}

\begin{theorem}[Literal Legendre synthesis and Jacobi closed form]
\label{thm:leg-mode-synthesis}
Let $\ell\ge0$ and let $M\in\PositiveIntegers$ satisfy
$\TwoAdicValuation(M)\ge\ell$.  Then, for $x\in[-1,1]$,
\begin{equation}
 \boxed{
 P_\ell(x)=\frac1M\sum_{|k|<2M}
 Q_\ell^-(k/M)u(x-k/M).}
 \label{eq:leg-mode-synthesis}
\end{equation}
The deconvolved polynomial has the finite Jacobi expansion
\begin{equation}
 \boxed{
 Q_\ell^-(x)=
 \sum_{r=0}^{\lfloor\ell/2\rfloor}
 \frac{\gamma_{2r}}{(2r)!}
 \frac{(\ell+2r)!}{2^{2r}\ell!}
 P_{\ell-2r}^{(2r,2r)}(x).}
 \label{eq:leg-Q-jacobi}
\end{equation}
In particular $Q_\ell^-(-x)=(-1)^\ell Q_\ell^-(x)$ and all its
coefficients are rational.
\end{theorem}

\begin{proof}
Equation \eqref{eq:leg-mode-synthesis} is
\cref{thm:lag-cardinal} applied to the polynomial $P_\ell$, or equivalently
\eqref{eq:interval-algorithm} with unit atom radius.  Expand
$\MM_u(D)^{-1}=\sum_r\gamma_rD^r/r!$.  Odd coefficients vanish by evenness,
and the classical derivative formula
\[
 D^{2r}P_\ell(x)=
 2^{-2r}\frac{(\ell+2r)!}{\ell!}
 P_{\ell-2r}^{(2r,2r)}(x)
\]
gives \eqref{eq:leg-Q-jacobi}.  Jacobi polynomials with integral parameters
have rational coefficients; every summand has the parity of $\ell$.
\end{proof}

% The next four subsections depend conceptually on the translated analysis
% kernel and biorthogonality proved later in this chapter.  Capture them here
% and typeset them only after that foundation, so the source stays close to
% its provenance while the printed exposition follows the proof dependency.
\long\def\LegendreDeferredClosure{%
\subsection{Reverse local closure and rational spectral identities}

For $x\in[-1,1]$ and $c\in[-2,2]$ put
\[
 I(c)=[-1,1]\cap[c-1,c+1].
\]
The shifted canonical expansion is
\begin{equation}
 u(x-c)=\mathbf1_{I(c)}(x)
 \sum_{r=0}^{\infty}a_rP_{2r}(x-c).
 \label{eq:leg-shifted-local}
\end{equation}
It is absolutely and uniformly convergent on
$[-1,1]\times[-2,2]$ because
$\sum_r|a_r|<\infty$ and $|P_{2r}|\le1$ on $[-1,1]$.
For $\ell,r\ge0$ and $M\in\PositiveIntegers$ with $\TwoAdicValuation(M)\ge\ell$, define
\begin{equation}
 H_{\ell r}^{(M)}(x)
 :=\frac1M\sum_{|k|<2M}Q_\ell^-(k/M)
 \mathbf1_{I(k/M)}(x)P_{2r}(x-k/M).
 \label{eq:leg-H}
\end{equation}
For $m,r\ge0$, define
\begin{equation}
 \mathscr C_{m,r}(c)
 :=\frac{2m+1}{2}\int_{I(c)}
 P_m(x)P_{2r}(x-c)\Differential x.
 \label{eq:leg-clipped-C}
\end{equation}

\begin{theorem}[Reverse local loop and rational spectral closure]
\label{thm:leg-reverse-loop}
Let $\ell\ge0$ and $M\in\PositiveIntegers$ satisfy $\TwoAdicValuation(M)\ge\ell$.  Then
\begin{equation}
 \boxed{P_\ell(x)=\sum_{r=0}^{\infty}a_rH_{\ell r}^{(M)}(x)}
 \label{eq:leg-reverse-loop}
\end{equation}
absolutely and uniformly on $[-1,1]$.  Moreover, for every $m\ge0$,
\begin{equation}
 \Lambda_m(c)=\sum_{r=0}^{\infty}a_r\mathscr C_{m,r}(c)
 \label{eq:leg-Lambda-connection}
\end{equation}
absolutely and uniformly for $c\in[-2,2]$.  For $m,r\ge0$ define
\begin{equation}
 T_{m\ell r}^{(M)}
 :=\frac1M\sum_{|k|<2M}
 Q_\ell^-(k/M)\mathscr C_{m,r}(k/M),
 \label{eq:leg-T}
\end{equation}
Then every $T_{m\ell r}^{(M)}$ is rational and
\begin{equation}
 \boxed{\delta_{m\ell}=\sum_{r=0}^{\infty}
 a_rT_{m\ell r}^{(M)}.}
 \label{eq:leg-rational-spectral}
\end{equation}
If $m+\ell$ is odd, then $T_{m\ell r}^{(M)}=0$ for every $r\ge0$.
\end{theorem}

\begin{proof}
Insert \eqref{eq:leg-shifted-local} into the finite synthesis
\eqref{eq:leg-mode-synthesis}.  Interchanging a finite sum with an absolutely
uniform series is legitimate and produces \eqref{eq:leg-reverse-loop}.
Insert the same expansion into \eqref{eq:leg-Lambda}; dominated uniform
convergence gives \eqref{eq:leg-Lambda-connection}.

Now substitute \eqref{eq:leg-Lambda-connection} into
\eqref{eq:leg-biorthogonality}.  The sum over $k$ is finite, so its
interchange with the absolutely convergent $r$-series gives
\eqref{eq:leg-rational-spectral}.  At rational $k/M$, the polynomial
$Q_\ell^-$ is rational, while \eqref{eq:leg-clipped-C} is the integral of a
rational polynomial over an interval with rational endpoints; hence every
$T$ is rational.  Finally,
\[
 Q_\ell^-(-c)\mathscr C_{m,r}(-c)
 =(-1)^{m+\ell}Q_\ell^-(c)\mathscr C_{m,r}(c),
\]
so odd total parity cancels in the symmetric $k$-sum.
\end{proof}

The functions $H_{\ell r}^{(M)}$ are only piecewise polynomial, but their
infinite weighted sum is the smooth polynomial $P_\ell$.  The cancellation of
all boundary jets is made explicit below.

\subsection{Gram geometry and finite projector frames}

Fix an integer $N\ge0$, put $M=4^N$, and write
$K_M:=\{k\in\IntegerNumbers:|k|<2M\}$.  Let
\[
 \phi_k(x)=u(x-k/M),\qquad
 \mathbf D_{kn}=M^{-1}Q_{2n}^-(k/M),\quad 0\le n\le N,
\]
and define
\[
 G_{k\ell}=\int_{-1}^{1}\phi_k(x)\phi_\ell(x)\Differential x,\qquad
 H=\operatorname{diag}(h_0,h_2,\ldots,h_{2N}).
\]

\begin{theorem}[Exact Gram isometry and coefficient projector]
\label{thm:leg-Gram}
The synthesis matrix satisfies
\begin{equation}
 \boxed{\mathbf D^TG\mathbf D=H.}
 \label{eq:leg-DtGD}
\end{equation}
Thus
\begin{equation}
 \Pi_{\rm coef}:=\mathbf D H^{-1}\mathbf D^TG
 \label{eq:leg-coef-projector}
\end{equation}
is an idempotent, is self-adjoint for the Gram seminorm
($G\Pi_{\rm coef}=\Pi_{\rm coef}^TG$), has range
$\operatorname{col}\mathbf D$, and has trace $N+1$.

The even Christoffel--Darboux kernel
\begin{equation}
 K_N^{\rm e}(x,y)=
 \sum_{n=0}^{N}\frac{P_{2n}(x)P_{2n}(y)}{h_{2n}}
 \label{eq:leg-even-CD}
\end{equation}
has the two-sided atom factorization
\begin{equation}
 \boxed{
 K_N^{\rm e}(x,y)=\frac1{M^2}
 \sum_{k,\ell\in K_M}
 \widetilde K_N^{\rm e}(k/M,\ell/M)
 u(x-k/M)u(y-\ell/M),}
 \label{eq:leg-two-sided-CD}
\end{equation}
where
\[
 \widetilde K_N^{\rm e}(s,t)
 =\sum_{n=0}^{N}\frac{Q_{2n}^-(s)Q_{2n}^-(t)}{h_{2n}}.
\]
Its coefficient matrix is positive semidefinite of rank $N+1$.
\end{theorem}

\begin{proof}
The $(m,n)$ entry of $\mathbf D^TG\mathbf D$ is the inner product of the finite syntheses of
$P_{2m}$ and $P_{2n}$, hence is $h_{2n}\delta_{mn}$.  The projector
identities follow by substituting $\mathbf D^TG\mathbf D=H$; cyclicity of trace gives
$\operatorname{tr}\Pi_{\rm coef}=N+1$.

Substitute \eqref{eq:leg-mode-synthesis} in both variables of
\eqref{eq:leg-even-CD} and interchange finite sums.  The coefficient matrix
is $\mathbf D H^{-1}\mathbf D^T$, hence positive semidefinite;
$\mathbf D$ has full column rank, so its
rank is $N+1$.
\end{proof}

There is also a one-sided synthesis/analysis factorization.  Put
\[
 K_N^\sharp(t,y)
 =\sum_{n=0}^{N}\frac{Q_{2n}^-(t)P_{2n}(y)}{h_{2n}},
\quad
 g_k(x)=u(x-k/M),\quad
 \eta_k(y)=M^{-1}K_N^\sharp(k/M,y).
\]
Then
\[
 \Pi_N^{\rm e}\colon L^2([-1,1])\longrightarrow
 \operatorname{span}\{P_0,P_2,\ldots,P_{2N}\}
\]
denotes the orthogonal projector, and
\begin{equation}
 K_N^{\rm e}(x,y)=\sum_{k\in K_M}g_k(x)\eta_k(y),
 \qquad
 (\Pi_N^{\rm e}f)(x)=\sum_{k\in K_M}g_k(x)\langle f,\eta_k\rangle.
 \label{eq:leg-one-sided-frame}
\end{equation}
It follows by applying polynomial synthesis in the $x$ variable.  Since
$\Pi_N^{\rm e}$ is an orthogonal projector of rank $N+1$,
\begin{align}
 \sum_k\langle g_k,\eta_k\rangle&=N+1,
 \label{eq:leg-frame-trace}\\
 \sum_{k,\ell}\langle g_k,g_\ell\rangle
 \langle\eta_k,\eta_\ell\rangle&=N+1.
 \label{eq:leg-frame-HS}
\end{align}
The first is the trace and the second the squared Hilbert--Schmidt norm of
the same projector.  These identities constrain any attempt to choose a
better-conditioned nullspace gauge.

\subsection{Refinement null trains and a lifting law}

For $p\in\Pi_d$ put $p^\sharp=\mathscr D_up$.  At a mesh
$m\in\PositiveIntegers$ with $\TwoAdicValuation(m)\ge d$, define the canonical finite coefficient
row
\[
 (F_mp)_k=m^{-1}p^\sharp(k/m),\qquad |k|<2m.
\]
Let $E_r$ embed a coarse row into the $r$-times finer grid by putting it at
indices divisible by $r$ and zero elsewhere.

\begin{theorem}[Interlevel null train and refinement cocycle]
\label{thm:leg-interlevel}
Let $d\ge0$, $p\in\Pi_d$, and $r,s,m\in\PositiveIntegers$ with
$\TwoAdicValuation(m)\ge d$.  Then, on $[-1,1]$,
\begin{equation}
 \boxed{
 \frac1{rm}\sum_{|k|<2rm}
 \bigl(1-r\mathbf1_{r\mid k}\bigr)
 p^\sharp(k/(rm))u(x-k/(rm))=0.}
 \label{eq:leg-interlevel}
\end{equation}
Equivalently, $G_{m;r}p:=F_{rm}p-E_rF_mp$ belongs to the local synthesis
kernel.  These gauge corrections satisfy
\begin{equation}
 \boxed{G_{m;rs}p=G_{rm;s}p+E_sG_{m;r}p.}
 \label{eq:leg-refinement-cocycle}
\end{equation}
\end{theorem}

\begin{proof}
Write the fine synthesis of $p$.  Rewrite the coarse synthesis on the fine
index set by setting $k=rj$; its coefficient becomes
$r(rm)^{-1}p^\sharp(k/(rm))$ on divisible indices and zero elsewhere.
Subtracting proves \eqref{eq:leg-interlevel}.  Both sides of
\eqref{eq:leg-refinement-cocycle} telescope to
$F_{rsm}p-E_sE_rF_mp$.
\end{proof}

For every integer $N\ge0$, let $a^{(N)}$ be the common-mesh coefficient row
for $S_N$:
\[
 a_k^{(N)}=4^{-N}C_N^-(k/4^N).
\]

\begin{corollary}[Exact quarter-grid Legendre lifting]
\label{cor:leg-lifting}
For every integer $N\ge0$, put $m=4^N$.  For $|k|<8m$,
\begin{equation}
 \boxed{
 a_k^{(N+1)}-(E_4a^{(N)})_k
 =\frac{1-4\mathbf1_{4\mid k}}{4m}C_N^-(k/(4m))
 +\frac{a_{N+1}}{4m}Q_{2N+2}^-(k/(4m)).}
 \label{eq:leg-lifting}
\end{equation}
The first term synthesizes zero on $[-1,1]$; the second synthesizes the new
orthogonal block $B_{N+1}$.
\end{corollary}

\begin{proof}
At $t=k/(4m)$,
\[
 a_k^{(N+1)}
 =\frac1{4m}\bigl(C_N^-(t)+a_{N+1}Q_{2N+2}^-(t)\bigr),
\]
whereas $(E_4a^{(N)})_k=
4\mathbf1_{4\mid k}C_N^-(t)/(4m)$.  Subtract.  Apply
\cref{thm:leg-interlevel} to $p=S_N$ for the first term and
\eqref{eq:leg-mode-synthesis} to the new mode for the second.
\end{proof}

\subsection{Flatness and finite boundary identities}

\begin{theorem}[Endpoint flatness and finite jet sum rules]
\label{thm:leg-flatness}
For every $s\ge0$,
\begin{equation}
 \boxed{
 \sum_{n\ge\lceil s/2\rceil}
 a_n\frac{(2n+s)!}{2^ss!(2n-s)!}=0,}
 \label{eq:leg-flatness-series}
\end{equation}
and the series is absolutely convergent.

For every $\ell\ge0$, every $M\in\PositiveIntegers$ with $\TwoAdicValuation(M)\ge\ell$, and
every $s\ge0$,
\begin{equation}
 \boxed{
 \frac1M\sum_{|k|<2M}Q_\ell^-(k/M)u^{(s)}(1-k/M)
 =
 \begin{cases}
 \dfrac{(\ell+s)!}{2^ss!(\ell-s)!},&s\le\ell,\\[0.7em]
 0,&s>\ell.
 \end{cases}}
 \label{eq:leg-boundary-jets}
\end{equation}
At the left endpoint the right side is multiplied by
$(-1)^{\ell+s}$.
\end{theorem}

\begin{proof}
Repeated self-adjointness of $\mathcal L$ shows that the Legendre
coefficients of the $C^\infty$ flat function $u$ decay faster than every
power.  Hence \eqref{eq:leg-canonical-series} may be differentiated any fixed
number of times.  Evaluate its $s$th derivative at $x=1$, use
$u^{(s)}(1)=0$, and use the classical formula
\[
 P_j^{(s)}(1)=
 \begin{cases}
 2^{-s}(j+s)!/[s!(j-s)!],&s\le j,\\
 0,&s>j.
 \end{cases}
\]
This proves \eqref{eq:leg-flatness-series} and its absolute convergence.

Differentiate the finite identity \eqref{eq:leg-mode-synthesis}; no
convergence argument is needed.  Evaluation at $1$ gives
\eqref{eq:leg-boundary-jets}, and reflection gives the left endpoint.
\end{proof}

\begin{corollary}[Boundary jets as finite Thue--Morse identities]
\label{cor:leg-thue-morse-boundary-jets}
Let $\ell,L,s$ be nonnegative integers with $L\ge\ell$, and set $M=2^L$.
For
$1\le k\le2M-1$, let $j_{k,s}\in\{0,\ldots,2^s-1\}$ and
$y_{k,s}\in[-1,1)$ be the unique cell coordinates satisfying
\begin{equation}
 1-\frac{k}{M}
 =-1+\frac{2j_{k,s}+1+y_{k,s}}{2^s}.
 \label{eq:leg-boundary-cell-coordinates}
\end{equation}
Then
\begin{equation}
 \boxed{
 \frac{2^{\binom{s+1}{2}}}{M}
 \sum_{k=1}^{2M-1}Q_\ell^-(k/M)
 (-1)^{\operatorname{wt}_2(j_{k,s})}u(y_{k,s})
 =
 \begin{cases}
 \dfrac{(\ell+s)!}{2^ss!(\ell-s)!},&s\le\ell,\\[0.7em]
 0,&s>\ell.
 \end{cases}}
 \label{eq:leg-thue-morse-boundary-jets}
\end{equation}
All values $y_{k,s}$ are dyadic.  Thus every displayed identity is an exact
rational cancellation combining the Bell--Bernoulli polynomial samples
$Q_\ell^-(k/M)$ with a level-$s$ Thue--Morse sign row.
\end{corollary}

\begin{proof}
The exact derivative-cell law is
\[
 u^{(s)}\!\left(-1+\frac{2j+1+y}{2^s}\right)
 =2^{\binom{s+1}{2}}(-1)^{\operatorname{wt}_2(j)}u(y),
 \qquad 0\le j<2^s,\quad -1\le y\le1.
\]
It is the theorem \path{iteratedDeriv_rvachev_cell} in the primary Lean
development.  The half-open cells in
\eqref{eq:leg-boundary-cell-coordinates} partition $[-1,1)$, so the stated
coordinates exist uniquely.  In \eqref{eq:leg-boundary-jets}, every term with
$k<0$ is evaluated strictly to the right of the support, while $k=0$ is the
flat endpoint $1$; all those derivatives vanish.  Apply the cell law to the
remaining indices $1\le k<2M$ and substitute.  At a shared cell boundary the
other convention has $y=1$ instead of $y=-1$; both values of $u$ vanish, so
the identity is independent of the convention.  Dyadicity and rationality
follow from the definitions and the exact dyadic-value theorem.
\end{proof}

\begin{corollary}[Central-binomial dyadic cancellation]
\label{cor:leg-central-sum}
For $n\ge0$ and $M=4^n$,
\begin{equation}
 \boxed{
 Q_{2n}^-(0)+2\sum_{k=1}^{M-1}
 Q_{2n}^-(k/M)u(k/M)=(-1)^n\binom{2n}{n}.}
 \label{eq:leg-central-sum}
\end{equation}
\end{corollary}

\begin{proof}
Evaluate \eqref{eq:leg-mode-synthesis} for $P_{2n}$ at $x=0$.
The boundary terms vanish, and parity pairs $k$ with $-k$.  Multiply by
$M=4^n$ and use
$P_{2n}(0)=(-1)^n4^{-n}\binom{2n}{n}$.
\end{proof}
}

For every integer $n\ge0$, put $m_n=4^n$ and define its finite even-mode train
\begin{equation}
 B_n(x):=\frac{a_n}{m_n}
 \sum_{|k|<2m_n}Q_{2n}^-(k/m_n)u(x-k/m_n).
 \label{eq:leg-block}
\end{equation}
Then $B_n=a_nP_{2n}$ on the target interval.

\begin{theorem}[Rank-one atomized spectral details]
\label{thm:leg-rank-one-details}
For every integer $n\ge0$ and $f\in L^2([-1,1])$, define
\begin{equation}
 \Delta_nf:=\frac{4n+1}{2}\langle f,P_{2n}\rangle P_{2n}.
 \label{eq:leg-rank-one-detail}
\end{equation}
Then $\Delta_n$ is the self-adjoint orthogonal rank-one projector onto
$\operatorname{span}\{P_{2n}\}$, and it has the finite atomized form
\begin{equation}
 \boxed{
 (\Delta_nf)(x)=\frac{4n+1}{2m_n}\langle f,P_{2n}\rangle
 \sum_{|k|<2m_n}Q_{2n}^-(k/m_n)u(x-k/m_n)}
 \quad(x\in[-1,1]).
 \label{eq:leg-rank-one-atomized}
\end{equation}
For all $m,n\ge0$,
\begin{equation}
 \boxed{\Delta_n\Delta_m=\delta_{nm}\Delta_n,
 \qquad \Delta_nu=B_n.}
 \label{eq:leg-rank-one-calculus}
\end{equation}
Thus the block loop is exactly the orthogonal spectral resolution of $u$,
with each spectral detail itself realized by a finite up-train.
\end{theorem}

\begin{proof}
The Legendre norm is
$\langle P_{2n},P_{2n}\rangle=2/(4n+1)$, so
\eqref{eq:leg-rank-one-detail} is the usual formula for the orthogonal
projector onto a one-dimensional subspace.  Legendre orthogonality gives
$\Delta_n\Delta_m=\delta_{nm}\Delta_n$.  Substituting the finite synthesis
\eqref{eq:leg-mode-synthesis} with $M=m_n$ proves
\eqref{eq:leg-rank-one-atomized}.  Finally,
$\langle u,P_{2n}\rangle=2a_n/(4n+1)$ by \eqref{eq:leg-an}, whence
$\Delta_nu=a_nP_{2n}=B_n$.
\end{proof}

\begin{theorem}[Blockwise self-loop and common-mesh partial sums]
\label{thm:leg-block-loop}
For every $x\in[-1,1]$,
\begin{equation}
 \boxed{u(x)=\sum_{n=0}^{\infty}B_n(x),}
 \label{eq:leg-block-loop}
\end{equation}
where the series converges absolutely and uniformly when each finite train is
kept as one block.

For an integer cutoff $N\ge0$, put
\begin{equation}
 S_N=\sum_{n=0}^{N}a_nP_{2n},\qquad
 C_N^-=\mathscr D_uS_N=\sum_{n=0}^{N}a_nQ_{2n}^-.
 \label{eq:leg-SN-CN}
\end{equation}
Then $\deg S_N\le2N$, and the sufficient common mesh $M_N=4^N$ gives
\begin{equation}
 \boxed{
 S_N(x)=\frac1{M_N}\sum_{|k|<2M_N}
 C_N^-(k/M_N)u(x-k/M_N).}
 \label{eq:leg-common-mesh}
\end{equation}
The right side converges uniformly to $u$.  The mesh $4^N$ is least for
universal exactness on the whole space $\Pi_{2N}$; no target-specific
minimality for $S_N$ is asserted.
\end{theorem}

\begin{proof}
Substituting \eqref{eq:leg-mode-synthesis} into
\eqref{eq:leg-canonical-series} gives \eqref{eq:leg-block-loop}; its $n$th
block is exactly $a_nP_{2n}$.  Since
$\|P_j\|_{L^\infty[-1,1]}\le1$ and $\sum_n|a_n|<\infty$, the Weierstrass
test proves absolute uniform block convergence.

The degree inequality in \eqref{eq:leg-SN-CN} is unconditional; equality
would require $a_N\ne0$.  Since $\TwoAdicValuation(4^N)=2N$, applying
\eqref{eq:leg-mode-synthesis} to the single polynomial $S_N$ gives
\eqref{eq:leg-common-mesh}.  Its convergence is precisely
$S_N\to u$.  The universal least-mesh statement is
\path{isLeast_rvachevCombExactThrough_even} from
\path{CompositeMeshSharpness.lean}, with the claim boundary already stated
after \cref{thm:scale-classification}.
\end{proof}

\begin{warning}[Blockwise does not mean atomwise]
Finite regrouping within $B_n$ is always legal, and the outer series may be
rearranged as a series of blocks.  It does not follow that the individual
atoms can be flattened into an absolutely or unconditionally convergent
double series.  The repeated central coefficient in
 \cref{thm:leg2-central-growth} proves the contrary.
\end{warning}

\subsection{Orthogonal blocks and positive energy series}

\begin{theorem}[Orthogonal translate blocks and exact tails]
\label{thm:leg-block-orthogonality}
For $m,n\ge0$,
\begin{equation}
 \boxed{
 \int_{-1}^{1}B_n(x)B_m(x)\Differential x
 =\delta_{nm}\frac{2a_n^2}{4n+1}.}
 \label{eq:leg-block-orthogonality}
\end{equation}
Consequently
\begin{align}
 \|u-S_N\|_2^2
 &=\sum_{n>N}\frac{2a_n^2}{4n+1},
 \label{eq:leg-L2-tail}\\
 \|u\|_2^2
 &=\sum_{n=0}^{\infty}\frac{2a_n^2}{4n+1}.
 \label{eq:leg-parseval}
\end{align}
For the Legendre Sturm--Liouville operator
\[
 \mathcal L=-D(1-x^2)D,\qquad
 \lambda_n=2n(2n+1),
\]
one further has
\begin{equation}
 \int_{-1}^{1}(1-x^2)B_n'B_m'\Differential x
 =\delta_{nm}\lambda_n\frac{2a_n^2}{4n+1},
 \label{eq:leg-derivative-blocks}
\end{equation}
and, for every real $s\ge0$ in the spectral domain,
\begin{equation}
 \|(I+\mathcal L)^{s/2}u\|_2^2
 =\sum_{n=0}^{\infty}(1+\lambda_n)^s
 \frac{2a_n^2}{4n+1}.
 \label{eq:leg-sobolev-energy}
\end{equation}
\end{theorem}

\begin{proof}
On $[-1,1]$, $B_n=a_nP_{2n}$, so
\eqref{eq:leg-block-orthogonality} is ordinary Legendre orthogonality.
Orthogonal convergence gives \eqref{eq:leg-L2-tail} and
\eqref{eq:leg-parseval}.  Since
$\mathcal LP_{2n}=\lambda_nP_{2n}$, integration by parts gives
\[
 \int(1-x^2)B_n'B_m'=\langle\mathcal LB_n,B_m\rangle,
\]
which proves \eqref{eq:leg-derivative-blocks}.  The last identity is the
spectral theorem for the nonnegative self-adjoint Legendre operator.
\end{proof}

The individual translates inside the blocks are not orthogonal.  Expanding
\eqref{eq:leg-block-orthogonality} instead gives an exact cross-scale
atom-Gram cancellation identity; this is the literal content formalized by
\path{sum_rvachevLegendreAtomCoefficient_mul_gram}.

Let
\begin{equation}
 \mathscr E_F:=\int_0^1\FabiusBounded(t)^2\Differential t.
 \label{eq:leg-A2}
\end{equation}
Because $u(x)=\FabiusBounded(1-|x|)$ on $[-1,1]$, $\|u\|_2^2=2\mathscr E_F$.

\begin{theorem}[Fabius square energy: rational series and sinc integrals]
\label{thm:leg-A2}
One has the three exact representations
\begin{align}
 \boxed{\mathscr E_F=\sum_{n=0}^{\infty}\frac{a_n^2}{4n+1}}
 \label{eq:leg-A2-series}\\
 &=\int_0^\infty
 \prod_{j=0}^{\infty}\SincPi^2(\xi/2^j)\Differential\xi
 \label{eq:leg-A2-fourier}\\
 &=\frac1{2\pi}\int_0^\infty
 \prod_{r=1}^{\infty}
 \left(\frac{\sin(t/2^r)}{t/2^r}\right)^2\Differential t.
 \label{eq:leg-A2-sinc}
\end{align}
Every finite partial sum in \eqref{eq:leg-A2-series} is a nonnegative
rational number, and the sequence of partial sums is nondecreasing.  It is
strict at the $n$th step exactly when $a_n\ne0$.
\end{theorem}

\begin{proof}
Divide \eqref{eq:leg-parseval} by two and use
$\|u\|_2^2=2\mathscr E_F$.  Rationality follows from $a_n\in\RationalNumbers$.
Parseval for the Fourier convention of \cref{sec:shared} gives
$\|u\|_2^2=\int_\RealNumbers|\widehat u(\xi)|^2\Differential\xi$.  Evenness, division by two,
and the sinc product give \eqref{eq:leg-A2-fourier}; the substitution
$t=2\pi\xi$ gives \eqref{eq:leg-A2-sinc}.  Monotonicity and its equality
criterion follow because the added term is $a_n^2/(4n+1)$.
\end{proof}

\begin{corollary}[Nested dyadic-Fabius square series]
\label{cor:leg-dyadic-square-energy}
For every integer $n\ge0$, define the finite dyadic quantity
\begin{equation}
 \mathcal R_n^{\rm dy}:=
 \sum_{j=0}^{n}(-1)^{n+j}
 \binom{2n}{n+j}\binom{2n+2j}{2n}(2j)!
 2^{\binom{2j+1}{2}}\FabiusBounded(2^{-2j-1}).
 \label{eq:leg-dyadic-R}
\end{equation}
Then, for every $n\ge0$, the coefficient identity in the following box
holds, and its substitution into \eqref{eq:leg-A2-series} gives the displayed
energy series:
\begin{equation}
 \boxed{a_n=4^{-n}(4n+1)\mathcal R_n^{\rm dy},
 \qquad
 \mathscr E_F=\sum_{n=0}^{\infty}
 (4n+1)16^{-n}(\mathcal R_n^{\rm dy})^2.}
 \label{eq:leg-dyadic-square-energy}
\end{equation}
Every inner quantity $\mathcal R_n^{\rm dy}$ is rational.  Hence this is a
positive outer series whose $n$th term is the square of a finite rational
combination of the nested values
$\FabiusBounded(2^{-1}),\FabiusBounded(2^{-3}),\ldots,\FabiusBounded(2^{-2n-1})$.
\end{corollary}

\begin{proof}
The coefficient identity in \eqref{eq:leg-dyadic-square-energy} is
kernel-verified as
\path{rvachevLegendreCoefficient_eq_fabius_sum} in
\path{FabiusLegendreCoefficients.lean}.  Here is its finite human-readable
derivation.
Put $\mu_{2j}=\int_{-1}^{1}x^{2j}u(x)\Differential x$.  Expanding $P_{2n}$ in
monomials in \eqref{eq:leg-an} gives
\[
 a_n=\frac{4n+1}{2^{2n+1}}
 \sum_{j=0}^{n}(-1)^{n+j}
 \binom{2n}{n+j}\binom{2n+2j}{2n}\mu_{2j}.
\]
The exact moment--Fabius bridge from the primary exposition is
\[
 \mu_{2j}=2(2j)!\,2^{\binom{2j+1}{2}}\FabiusBounded(2^{-2j-1});
\]
obtained by integrating the dyadic derivative-cell formula; its dyadic
values are kernel-verified by \path{fabiusDyadic_cast}.  Substitution proves
the formula for $a_n$.  Insert its square into
\eqref{eq:leg-A2-series} and cancel one factor $4n+1$ to obtain the second
identity.  The same exact dyadic-value theorem gives rationality.
\end{proof}

The four exact cutoffs
\[
 \frac14,\qquad\frac7{18},\qquad
 \frac{3271}{8100},\qquad
 \frac{3246043}{8037225}
\]
are kernel-verified in \path{FabiusLegendreRationalEnergy.lean}; the Fourier
and sinc equalities are kernel-verified in
\path{FabiusSquareEnergyFourier.lean}.  A decimal evaluation of
$\mathscr E_F$ is not
needed for any theorem here.

\subsection{Translated Legendre analysis and exact biorthogonality}

\begin{definition}[Translated Legendre analysis kernel]
\label{def:leg-Lambda}
For $m\ge0$ and $c\in\RealNumbers$, set
\begin{equation}
 \Lambda_m(c)=\frac{2m+1}{2}
 \int_{-1}^{1}u(x-c)P_m(x)\Differential x.
 \label{eq:leg-Lambda}
\end{equation}
\end{definition}

\begin{theorem}[Analysis kernel and Fourier--Bessel representation]
\label{thm:leg-Lambda}
For every $m\ge0$, the function $\Lambda_m$ belongs to $C_c^\infty(\RealNumbers)$,
has support contained
in $[-2,2]$, and satisfies
\begin{align}
 \Lambda_m(-c)&=(-1)^m\Lambda_m(c),
 \label{eq:leg-Lambda-parity}\\
 \Lambda_{2n}(0)&=a_n,\qquad \Lambda_{2n+1}(0)=0\quad(n\ge0),
 \label{eq:leg-Lambda-origin}\\
 \Lambda_m(c)&=(2m+1)(-\ImaginaryUnit)^m
 \int_\RealNumbers\widehat u(\xi)j_m(2\pi\xi)
 \EulerE^{2\pi\ImaginaryUnit\xi c}\Differential\xi.
 \label{eq:leg-Lambda-bessel}
\end{align}
Here $j_m$ is the spherical Bessel function.
\end{theorem}

\begin{proof}
Let $\chi_m(x)=P_m(x)\mathbf1_{[-1,1]}(x)$.  Evenness of $u$ gives
$\Lambda_m=(2m+1)(u*\chi_m)/2$.  Convolution of a compactly supported
$C^\infty$ function with a compactly supported integrable function is
$C^\infty$, and supports add to $[-2,2]$.  Reflection gives
\eqref{eq:leg-Lambda-parity}; \eqref{eq:leg-Lambda-origin} is
\eqref{eq:leg-an} and oddness.

The classical integral representation gives
$\widehat\chi_m(\xi)=2(-\ImaginaryUnit)^mj_m(2\pi\xi)$.  The rapid decay of
$\widehat u$ makes the product integrable, so Fourier inversion of the
convolution proves \eqref{eq:leg-Lambda-bessel}.
\end{proof}

\begin{lemma}[Rational dyadic truncated moments]
\label{lem:leg-dyadic-truncated}
If $a,b\in\IntegerNumbers[1/2]$, $a\le b$, and $r\ge0$, then
\begin{equation}
 \int_a^b x^ru(x)\Differential x\in\RationalNumbers.
 \label{eq:leg-dyadic-truncated}
\end{equation}
Consequently $\Lambda_m(c)\in\RationalNumbers$ for every dyadic $c$.
\end{lemma}

\begin{proof}
Intersect $[a,b]$ with the support $[-1,1]$, which preserves dyadic
endpoints.  On $[-1,1]$ the differential refinement equation implies
\[
 u(x)=\frac12u'\!\left(\frac{x-1}{2}\right),
\]
because the other child $u(x-2)$ vanishes.  For a polynomial $p$ and dyadic
$a,b\in[-1,1]$, substitute $t=(x-1)/2$ and integrate by parts:
\begin{align*}
 \int_a^bp(x)u(x)\Differential x
 &=
 \left[p(2t+1)u(t)\right]_{(a-1)/2}^{(b-1)/2}\\
 &\quad-\int_{(a-1)/2}^{(b-1)/2}
 2p'(2t+1)u(t)\Differential t.
\end{align*}
The boundary values of $u$ at dyadic arguments are rational by
\path{fabiusDyadic_cast} in \path{DyadicAnalytic.lean}, the exact
dyadic-value theorem recorded in the primary exposition.  Induction on $\deg p$
therefore proves the first claim, with the constant case terminating after
the boundary term.

For dyadic $c$, substitute $x=c+y$ in \eqref{eq:leg-Lambda}.  Expand
$P_m(y+c)$ and split at the dyadic endpoints of
$[-1,1]\cap[-1-c,1-c]$.  The result is a finite rational combination of the
moments just proved.
\end{proof}

\begin{theorem}[Legendre--up dyadic biorthogonality]
\label{thm:leg-biorthogonality}
Let $\ell\ge0$, $M\in\PositiveIntegers$ with $\TwoAdicValuation(M)\ge\ell$, and $m\ge0$.
Then
\begin{equation}
 \boxed{
 \frac1M\sum_{|k|<2M}
 Q_\ell^-(k/M)\Lambda_m(k/M)=\delta_{m\ell}.}
 \label{eq:leg-biorthogonality}
\end{equation}
\end{theorem}

\begin{proof}
Multiply \eqref{eq:leg-mode-synthesis} by
$(2m+1)P_m(x)/2$ and integrate over $[-1,1]$.  The left side is
$\delta_{m\ell}$ by Legendre orthogonality, and each atom integral is
$\Lambda_m(k/M)$.
\end{proof}

\begin{corollary}[Finite biorthogonal matrices and projector identities]
\label{cor:leg-biorthogonal-matrices}
Let $d\ge0$ and $M\in\PositiveIntegers$ satisfy $\TwoAdicValuation(M)\ge d$.  Then the matrices
indexed by $K_M:=\{k\in\IntegerNumbers:|k|<2M\}$,
\[
 A_{k\ell}^{(d)}=M^{-1}Q_\ell^-(k/M),\qquad
 L_{mk}^{(d)}=\Lambda_m(k/M)
 \quad(k\in K_M,\ 0\le m,\ell\le d)
\]
satisfy $L^{(d)}A^{(d)}=I_{d+1}$.  Consequently
$A^{(d)}L^{(d)}$ is an oblique projector of rank and trace $d+1$, and
\begin{align}
 d+1&=\frac1M\sum_{|k|<2M}\sum_{\ell=0}^{d}
 Q_\ell^-(k/M)\Lambda_\ell(k/M),
 \label{eq:leg-spectral-trace}\\
 1&=\sum_{\substack{S\subset K_M\\|S|=d+1}}
 \det L^{(d)}_{:,S}\det A^{(d)}_{S,:}.
 \label{eq:leg-spectral-CB}
\end{align}
\end{corollary}

\begin{proof}
Because $\ell\le d$, the hypothesis $\TwoAdicValuation(M)\ge d$ permits
\cref{thm:leg-biorthogonality} for every matrix entry, so
$L^{(d)}A^{(d)}=I_{d+1}$.  Idempotence, rank, and trace follow from this
right-inverse identity;
taking the trace and applying Cauchy--Binet give
\eqref{eq:leg-spectral-trace} and \eqref{eq:leg-spectral-CB}.
\end{proof}

\LegendreDeferredClosure
